% !TEX root = ../AImath.v5.tex

\chapter{ 前言}
\section{从"困难重重"到"显而易见"：智能发展的节奏}

% ChatGPT建议：统一为课堂讲解语气；段首补承接提示，段尾加入启发句；去除同义重复；保持自然衔接与LaTeX格式。
% 修改后版本
回顾人工智能的发展，我们反复看到一种熟悉的节奏：昨天还被认为“困难重重”的问题，今天却变得“顺理成章”。从早期的专家系统到当下的大语言模型，从手工特征工程到端到端的深度学习，这条演进主线从未中断，只是以新的形式不断重现。许多被视为“难题”的问题，往往并非不可解，而是解决方法在当时尚未成熟。

这种由“困难”走向“显而易见”的转变并非偶然，而是研究范式成熟的结果。当我们抓住了问题的本质，并匹配到恰当的数学工具与计算方法，看似不可逾越的障碍就会自然消解。正如牛顿所言：“如果我看得更远，那是因为我站在巨人的肩膀上。”每一次跨越，都是在前人成果之上完成的再组织与再表达。从这个意义上说，我们今天所认为的“显而易见”，或许正孕育着明天的新突破。

\section{数学思维：智能发展的主线}
% ChatGPT建议：在段首增加思维引导句；补充过渡句连接三段；保持“数学—智能”主线；调整语气为启发式课堂体。
% 修改后版本
要理解人工智能的根脉，我们几乎绕不开数学。数学不仅是描述世界的语言，更是推动智能演化的主线。从古希腊的几何学到现代的概率论，从牛顿的微积分到香农的信息论，每一个数学概念的诞生，都为“理解”与“计算”智能提供了新的可能。

%-----这一段待商榷和修改-------
线性代数让我们得以处理高维数据，微积分使得梯度下降成为可能，概率论为不确定性建模奠定了基础，而信息论则揭示了学习与压缩的内在联系。回望人工智能的发展，每一次重大突破几乎都伴随着数学思想的革新或重生。
每一次理论的成熟，都在为“理解智能”铺路。
%----

深度学习的成功，本质上是这些数学思想的再整合。
%-----这一段待商榷和修改-------
神经网络的反向传播算法体现了微积分的链式法则，卷积神经网络继承了信号处理的卷积理论，而注意力机制则借鉴了概率分布与加权平均的思维。
这一切都说明：数学不仅是工具，更是一种理解智能的语言。

\section{智能的层次：从符号到学习，再到创造}
回望人工智能的发展史，我们可以看到一条清晰的演化脉络：从以逻辑推理为代表的“符号主义智能”，到以数据驱动为核心的“学习型智能”，再到当下正逐渐显现的“创造性智能”。这不仅是技术范式的更迭，更是人类理解“智能”的方式在不断深化。

最初阶段的人工智能以符号与规则推理为基础。二十世纪中叶的第一代研究者相信，只要能把人类的知识形式化，机器就能像人类一样思考。早期的专家系统与定理证明程序正是这一理念的体现——它们擅长“演绎”，却难以处理模糊性与经验的不确定。

随着数据与计算能力的增长，人工智能进入了“学习”的时代。机器不再依赖人工制定的规则，而是从海量数据中自动发现模式与规律。这一阶段的智能更接近“模仿”：系统通过神经网络与统计模型，复现人类在语言、视觉和决策上的能力。正如孩子学说话，通过反复尝试与纠错，机器逐渐掌握了理解世界的路径。

当学习能力积累到一定程度，智能便开始显露“创造”的潜质。所谓创造，并非凭空发明，而是在已有知识与模式的基础上重新组合、抽象与泛化，从而生成新的结构与洞见。今天的大语言模型与生成式系统，让我们首次直观地感受到这种“新智能”的雏形：它们不仅能模仿，还能在语义与逻辑的空间中创造出人类未曾明确教给它们的内容。

从符号到学习，再到创造，人工智能的演化实际上映射了人类认知方式的演进——从理性逻辑到经验学习，再到综合创新。真正的“创造”，或许正是理解与模仿的高度统一。这一层层的演化，也构成了我们理解智能的路径。

\section{在浪潮中理解智能}
2025年的今天，几乎每隔几天就会诞生一个人工智能新工具。新的模型、框架和平台层出不穷。
%人们惊叹于AI的速度，也常感到跟不上它的脚步。
技术越快，人越焦虑。技术的指数级进步，让我们体验到前所未有的便利，也带来了同样前所未有的不安——焦虑于被取代、被落下，甚至焦虑于自己理解能力的局限。  

面对这股奔腾的技术浪潮，真正值得思考的问题是：在这一切变化背后，人工智能的本质究竟是什么？我们该如何在更新速度如此惊人的时代，建立一种不被浪潮裹挟的理性理解？  

本书的出发点，是希望与读者一道，在纷繁的技术表象之下，寻找理解人工智能的更深层逻辑。我们并不满足于对算法的罗列与介绍，而是想通过这本书的尝试，揭示其背后的数学结构、历史脉络与思想根基。唯有如此，我们才能真正“看懂”智能的成长线索，而不是被无尽的更新所淹没。  

{\bf 学习的关键，不在于记忆全部，而在于洞见规律。}

正如《一如既往》一书所指出的：“
{\bf 刚开始研究一个领域时，往往感觉有无数的知识点需要记忆。其实并非如此。}
如果能够识别统领这个领域的核心法则——通常是3$\sim$12条——
那么当初看似必须记住的数百万件事物，其实只是这些原则的不同组合。”
理解本质规律，就能在纷繁变化中保持清醒与方向。  

因此，我们并不是“讲授者”，而和大家同行的学习者。  
这本书既是对人工智能数学思维的探索，也是一次共同学习的过程。  
如果在阅读中，你能对某个概念产生新的联想，或在某个例子中看到自己的理解路径——那便是我们写作的意义所在。  

%最大的挑战在于：如何在保持严谨的同时，让内容易于理解。{\bf 数学是一种极为精确的语言，但若缺乏直觉引导，也容易成为理解的障碍。}为此，我们提出以下策略：
%\begin{itemize}
%\item \emph{从直觉出发}：以生活化或熟悉的例子引入，再逐步深入数学细节；
%\item \emph{历史视角}：通过历史演进追溯思想的形成；
%\item \emph{实践结合}：理论与应用相辅相成，让抽象变得可感；
%\item \emph{层次递进}：从简单到复杂，从具体到抽象，循序渐进地搭建知识体系。
%\end{itemize}

带着问题阅读，带着好奇探索——这正是学习的起点，也是智能成长的本源。

\newpage